\scanpage{377}
\mspage{116}
% archive sheet number: 376
\[
\tag*{$\unsure{(\alpha)|(2)}$}
H_T(a,b)\mathrel{\unsure{-|\to}}H_T(b,a),\quad c\mathrel{\unsure{\mapstochar\relbar|\mapsto}}c^\circ,
\]
donnant lieu au diagramme commutatif exprimant que c'est bien un inverse pour la composition. Une autre façon de le dire, c'est que pour \unsure{$\mathrm V{\Rightarrow}$|tout} objet $Y$ de $\mathrm{Cat\,hot}$ (ou de $\mathrm{Ind\,Cat\,hot}$), \deleted{\unsure{que}} les ensembles
\[
\mathrm{Hom}_{\mathrm{Cat\,hot}}(Y,H_T(a,b))\overset{\text{\unsure{def|déf}}}{=}H_T^Y(a,b)
\]
forment, pour $a,b\in X$ variables, \inserted{les flèches d'}un \uline{groupoïde}$^{*}$ \inserted{$\mathrm{Ch}_{\Sigma}^{Y}(X)$} ayant pour \inserted{ensemble de} sommets $\mathrm{Ob}\,X$, et $H_T^Y(a,b)$ étant formé des flèches de $\unsure{a_x|a}$ vers $b_{\deleted{\unsure{a}}}$, \unsure{et|les} \deleted{\unsure{fonct}} l'application
% slanted
\margin{Si $\Sigma=\mathrm{Fl}(X)$. Si $\Sigma$ est plus petit, au lieu d'un groupoïde, on trouve une \uline{catég\unsure{orie}} sans plus.}
\[
\deleted{\unsure{\mathrm{Ch}_{\Sigma}}}\ H_T^Y(a,b)\mathrel{-}H_T^Y(b,a),\qquad c\mathrel{\mapstochar\relbar}c^\circ
\]
\deleted{\unsure{fonct}} \unsure{déduite|déduit} du foncteur $(\alpha)$ plus haut n'est autre que l'application de passage à l'inverse. Cette \deleted{groupoïde} \inserted{catégorie} n'est d'ailleurs autre que celle déduite de
\[
X(Y)\Sigma_{\unsure{Y|y}}^{-1}\qquad\text{où}\quad
\left\{\begin{aligned}
&X(Y)=\underline{\mathrm{Hom}}_{\mathrm{Cat}}(Y,X)\\
&\Sigma_Y\subset\mathrm{Fl}(X_Y)\quad\text{défini}\\
&\qquad\text{en termes de }\Sigma\\
&\qquad\text{``argument par}\\
&\qquad\text{argument'',}
\end{aligned}\right.
\]
à l'aide du \deleted{\unsure{$\mathrm{fl.}$}} foncteur
\[
X\mathrel{-}X_Y,\qquad a\mathrel{\mapstochar\relbar}a_Y
\]
